%
% new_tokens_label_guard_buggy.tex
%
% A Z specification of the luxd-pty daemon's memory of retry tokens, the
% entry it keeps for each launch it admits from a `new` request's token.
% Fidelity variant: label_guard.
%
% Design record: DESIGN.md, DES-107 (token memory, "Formal model")
% Type-check:  fuzz -t docs/new_tokens_label_guard_buggy.tex
% Model-check: scripts/verify_new_tokens_model.sh
%
\documentclass[a4paper,10pt,fleqn]{article}
\usepackage[margin=1in]{geometry}
\usepackage{fuzz}

\begin{document}

\title{Retry-Token Memory, Fidelity Variant label guard}
\author{Formal model of the daemon's retry-token memory}
\date{November 2026}
\maketitle

\section{Introduction}

The retry guard is the label, which is freed on exit, so a launch that succeeded and exited launches again on retry. Fails T1.

Everything else is the intact specification, \texttt{new\_tokens.tex}; the variant is made to fail the invariants named below.

\section{Basic Types}

\begin{zed}
  [TOK] \\
  LAB ::= lbl1 | lbl2
\end{zed}

\begin{axdef}
  cap : \nat
\where
  cap = 2
\end{axdef}

\texttt{TOK} is a retry token.  \texttt{LAB} is a session label.  The
contract's 256 is scaled to $cap = 2$ with three tokens, which is enough to
evict one entry and to retry it.

\section{State}

\begin{schema}{Sys}
  held : TOK \pfun (0 \upto cap) \\
  bound, sv : \power TOK \\
  lab : LAB \pfun TOK \\
  lv, nforks : TOK \fun (0 \upto 2) \\
  recent : \seq TOK
\where
  bound \subseteq \dom held \\
  \# recent \leq cap
\end{schema}

\texttt{held} is the memory: a token maps to the age of its entry.
\texttt{bound} are the held tokens whose session was recorded.  \texttt{lab}
maps a label to the token of the live session that owns it.  \texttt{lv} counts
the live processes of a token.  \texttt{sv} are the tokens whose shell survived
a restart.  \texttt{recent} and \texttt{nforks} are the history variables.

\begin{schema}{Init}
  Sys'
\where
  held' = \emptyset \\
  bound' = \emptyset \\
  sv' = \emptyset \\
  lab' = \emptyset \\
  lv' = TOK \cross \{ 0 \} \\
  nforks' = TOK \cross \{ 0 \} \\
  recent' = \langle \rangle
\end{schema}

\section{Operations}

\begin{schema}{RepeatBound}
  \Delta Sys \\
  t? : TOK \\ l? : LAB
\where
  l? \in \dom lab \\
  lab~l? = t? \\
  held' = held \\
  bound' = bound \\
  lab' = lab \\
  lv' = lv \\
  sv' = sv \\
  recent' = recent \\
  nforks' = nforks
\end{schema}

A repeated token whose session was recorded is answered with the original session, running or exited. Nothing changes.

\begin{schema}{RefuseLabel}
  \Delta Sys \\
  t? : TOK \\ l? : LAB
\where
  \lnot (l? \in \dom lab \land lab~l? = t?) \\
  l? \in \dom lab \\
  held' = held \\
  bound' = bound \\
  lab' = lab \\
  lv' = lv \\
  sv' = sv \\
  recent' = recent \\
  nforks' = nforks
\end{schema}

A launch refused because its label is held by a live session. The refusal precedes the fork and adds no entry.

\begin{schema}{RefuseConfig}
  \Delta Sys \\
  t? : TOK \\ l? : LAB
\where
  \lnot (l? \in \dom lab \land lab~l? = t?) \\
  held' = held \\
  bound' = bound \\
  lab' = lab \\
  lv' = lv \\
  sv' = sv \\
  recent' = recent \\
  nforks' = nforks
\end{schema}

Any other refusal before the fork: runtime launchers off, unknown launcher. No entry is added.

\begin{schema}{LaunchOk}
  \Delta Sys \\
  t? : TOK \\ l? : LAB
\where
  \lnot (l? \in \dom lab \land lab~l? = t?) \\
  l? \notin \dom lab \\
  lv~t? < 2 \\
  nforks~t? < 2 \\
  held' = \{ x : TOK; a : 0 \upto cap | (x \mapsto a) \in (\{ t? \} \ndres held) \land a + 1 < cap @ x \mapsto a + 1 \} \cup \{ t? \mapsto 0 \} \\
  bound' = (bound \cap \dom held') \cup \{ t? \} \\
  lab' = lab \cup \{ l? \mapsto t? \} \\
  lv' = lv \oplus \{ t? \mapsto lv~t? + 1 \} \\
  sv' = sv \\
  ((\# recent < cap \land recent' = recent \cat \langle t? \rangle) \lor (\# recent = cap \land (\exists h : TOK; r : \seq TOK @ recent = \langle h \rangle \cat r \land recent' = r \cat \langle t? \rangle))) \\
  nforks' = (\ran recent' \dres (nforks \oplus \{ t? \mapsto nforks~t? + 1 \})) \cup ((TOK \setminus \ran recent') \cross \{ 0 \})
\end{schema}

An admitted launch whose fork and session both succeed. The entry is added before the fork and bound to the session.

\begin{schema}{LaunchForkFails}
  \Delta Sys \\
  t? : TOK \\ l? : LAB
\where
  \lnot (l? \in \dom lab \land lab~l? = t?) \\
  l? \notin \dom lab \\
  held' = held \\
  bound' = bound \\
  lab' = lab \\
  lv' = lv \\
  sv' = sv \\
  recent' = recent \\
  nforks' = nforks
\end{schema}

The fork fails. The entry added before the fork is removed and the oldest entry its addition dropped is restored, so the memory is as it was, and the reply is \texttt{fault}.

\begin{schema}{LaunchUnrecorded}
  \Delta Sys \\
  t? : TOK \\ l? : LAB
\where
  \lnot (l? \in \dom lab \land lab~l? = t?) \\
  l? \notin \dom lab \\
  lv~t? < 2 \\
  nforks~t? < 2 \\
  held' = \{ x : TOK; a : 0 \upto cap | (x \mapsto a) \in (\{ t? \} \ndres held) \land a + 1 < cap @ x \mapsto a + 1 \} \cup \{ t? \mapsto 0 \} \\
  bound' = bound \cap \dom held' \\
  lab' = lab \\
  lv' = lv \oplus \{ t? \mapsto lv~t? + 1 \} \\
  sv' = sv \\
  ((\# recent < cap \land recent' = recent \cat \langle t? \rangle) \lor (\# recent = cap \land (\exists h : TOK; r : \seq TOK @ recent = \langle h \rangle \cat r \land recent' = r \cat \langle t? \rangle))) \\
  nforks' = (\ran recent' \dres (nforks \oplus \{ t? \mapsto nforks~t? + 1 \})) \cup ((TOK \setminus \ran recent') \cross \{ 0 \})
\end{schema}

The entry is added, the process forks, and the session is never recorded.

\begin{schema}{ExitLabelled}
  \Delta Sys \\
  l? : LAB
\where
  l? \in \dom lab \\
  lv~(lab~l?) > 0 \\
  held' = held \\
  bound' = bound \\
  lab' = \{ l? \} \ndres lab \\
  lv' = lv \oplus \{ lab~l? \mapsto lv~(lab~l?) - 1 \} \\
  sv' = sv \\
  recent' = recent \\
  nforks' = nforks
\end{schema}

A recorded session exits and frees its label. Its token stays held and bound.

\begin{schema}{ExitStray}
  \Delta Sys \\
  t? : TOK
\where
  lv~t? > \# (lab \rres \{ t? \}) \\
  held' = held \\
  bound' = bound \\
  lab' = lab \\
  lv' = lv \oplus \{ t? \mapsto lv~t? - 1 \} \\
  sv' = sv \\
  recent' = recent \\
  nforks' = nforks
\end{schema}

A process with no recorded session exits.

\begin{schema}{RestartClean}
  \Delta Sys
\where
  held' = \emptyset \\
  bound' = \emptyset \\
  lab' = \emptyset \\
  lv' = (TOK \cross \{ 0 \}) \oplus \{ x : TOK | x \in sv \land lv~x > 0 @ x \mapsto 1 \} \\
  sv' = sv \\
  recent' = \langle \rangle \\
  nforks' = TOK \cross \{ 0 \}
\end{schema}

The daemon restarts. The memory is cleared and every shell the restart can end is ended. A shell that survived an earlier restart is out of its reach.

\begin{schema}{RestartSurvivor}
  \Delta Sys \\
  s? : \power TOK
\where
  s? \neq \emptyset \\
  \forall x : s? @ lv~x > 0 \\
  held' = \emptyset \\
  bound' = \emptyset \\
  lab' = \emptyset \\
  lv' = (TOK \cross \{ 0 \}) \oplus \{ x : TOK | x \in sv' \land lv~x > 0 @ x \mapsto 1 \} \\
  sv' = sv \cup s? \\
  recent' = \langle \rangle \\
  nforks' = TOK \cross \{ 0 \}
\end{schema}

The daemon restarts and a shell for each token in $s?$ survives SIGKILL. The daemon logs its token and process id and exits with a failure status; the survivor has no history.

\section{Invariants}

Each invariant is checked as a violation predicate in
\texttt{scripts/verify\_new\_tokens\_model.sh}.

\begin{description}
  \item[T1, a held token never yields a second launch.]
    $\forall t : TOK @ nforks~t \leq 1$.
  \item[T2, refusals and repeats add no entry.]
    $\dom held \subseteq \ran recent$: every entry belongs to an admitted
    launch of the last $cap$.
  \item[T3, the oldest is dropped only at the cap.]
    $\{ t : \ran recent | nforks~t > 0 \} \subseteq \dom held$: no token that
    forked in the last $cap$ admitted launches has been dropped early.  With T2
    and $\# recent \leq cap$ the memory is exactly the entries of the last
    $cap$ launches that forked.
  \item[T4, a launch with no process can be retried.]
    $\forall t : \dom held @ nforks~t > 0$: an entry exists only for a token
    whose launch forked.
  \item[T5, the memory is bounded.] $\# (\dom held) \leq cap$.
\end{description}


\end{document}
